% Figure: schematic engineering stress-strain curve for a ductile metal, with
% the elastic region, yield point, strain-hardening region, ultimate stress,
% necking, and fracture marked. Plus a brittle curve for contrast.
\begin{figure}[htbp]
\centering
\begin{tikzpicture}
\begin{axis}[
  width=12cm, height=8cm,
  xlabel={engineering strain $\epsilon$}, ylabel={engineering stress $\sigma$ (MPa)},
  xmin=0, xmax=0.30, ymin=0, ymax=520,
  axis lines=left,
  legend pos=south east,
  legend cell align=left,
  tick label style={font=\scriptsize},
  label style={font=\small},
  legend style={font=\scriptsize},
  clip=false,
]
% Ductile mild steel: linear to yield ~250 at eps~0.00125, yield plateau,
% strain hardening to UTS ~450 at eps~0.18, then necking down to fracture ~380 at 0.27.
\addplot[engblue, very thick] coordinates {
  (0,0) (0.00125,250)
};
\addplot[engblue, very thick] coordinates {
  (0.00125,250) (0.02,255)
};
\addplot[engblue, very thick, smooth] coordinates {
  (0.02,255) (0.05,330) (0.09,400) (0.14,440) (0.18,450)
};
\addplot[engblue, very thick, smooth] coordinates {
  (0.18,450) (0.22,430) (0.26,390) (0.27,375)
};
\addlegendentry{ductile metal (mild steel)}
% fracture marker
\node[engblue] at (axis cs:0.27,375) {\large $\times$};
% brittle material: nearly linear to a sudden fracture, no plateau.
\addplot[engred, very thick] coordinates {
  (0,0) (0.012,500)
};
\addlegendentry{brittle material (cast iron, schematic)}
\node[engred] at (axis cs:0.012,500) {\large $\times$};
% annotations
\node[engblue, font=\scriptsize, anchor=west] at (axis cs:0.003,250) {yield $\sigma_y$};
\node[engblue, font=\scriptsize, anchor=south] at (axis cs:0.18,455) {ultimate $\sigma_u$};
\node[engblue, font=\scriptsize, anchor=north west] at (axis cs:0.20,420) {necking};
\node[font=\scriptsize, anchor=north west] at (axis cs:0.001,120) {elastic, slope $E$};
\end{axis}
\end{tikzpicture}
\caption[Engineering stress-strain curves, ductile and brittle]{Engineering
stress-strain curves. The ductile metal (blue) rises linearly through the
elastic region with slope $E$, reaches the yield stress $\sigma_y$ where
irreversible deformation begins, passes through a short yield plateau and a
strain-hardening region up to the ultimate stress $\sigma_u$, then the
specimen necks and the engineering stress falls until fracture (the cross).
The brittle material (red) is nearly linear to a sudden fracture with little
preceding plastic deformation. Engineering stress uses the original
cross-sectional area, which is why the ductile curve turns down after
necking even though the true stress on the shrinking neck keeps rising.}
\label{fig:vol03ch03:stress-strain-curve}
\end{figure}
